%%%% CAPÍTULO 4: RESULTADOS PARCIAIS E DISCUSSÃO
%%
%% Deve descrever detalhadamente os dados obtidos no trabalho. Os resultados são
%% normalmente discutidos a partir de ilustrações (gráficos, quadros, etc.),
%% tabelas, entre outros elementos, que podem ser incluídos no documento. Deve
%% efetuar a comparação dos dados obtidos e/ou resultados com aqueles descritos
%% na revisão da literatura, incluindo os comentários sobre os estudos de outros
%% trabalhos.

%% Locais (pastas) de ilustrações deste capítulo
\graphicspath{%
  {./Chapter-4/}%% Primário
%   {./Chapter-4/Illustrations/}%% Secundário (descomentar se houver)
}

\chapter{Partial Results and Discussion}
\label{sec:results}
 
This section evaluates PID-transit on two datasets: the Porto metropolitan bus
feed introduced in Section~\ref{subsec:data_layer}, and the Swiss national NeTEx
timetable, a 24.5\,GB production dataset published by an independent authority.
Section~\ref{sec:res-ingestion} reports GTFS ingestion
and mapping to the Transmodel relational schema. Section~\ref{sec:res-validation}
reports validator outcomes for round-trip export to GTFS and NeTEx and quantifies
the warning-level changes induced by relational normalization.
Section~\ref{sec:res-performance} reports wall-clock duration, artifact sizes, and
the compression obtained by journey-pattern deduplication.
Section~\ref{sec:res-netex-scale} evaluates the NeTEx import path at national
scale on the Swiss feed, and Section~\ref{sec:res-benchmark} benchmarks
PID-transit against one representative of each competing architectural class.
 
\section{Network Ingestion and Relational Mapping}
\label{sec:res-ingestion}
 
The PID-transit toolkit is released as an open-source Python package under the MIT
license (available on request). GTFS entities
from the case-study feed (Section~\ref{sec:data-source}) were mapped to their
Transmodel equivalents and persisted in the relational database;
Table~\ref{tab:import_results} reports the resulting entity counts.
 
\begin{table}[htbp]
  \centering
  \caption{Transmodel entities created from the GTFS feed.}
  \label{tab:import_results}
  \begin{tabular}{l r}
    \hline
    \textbf{Transmodel Entity} & \textbf{Count} \\
    \hline
    Operators                    & 1        \\
    Feed Info                    & 1        \\
    Scheduled Stop Points        & 2\,550   \\
    Lines                        & 71       \\
    Routes                       & 138      \\
    Day Types                    & 3        \\
    Operating Day Exceptions     & 365      \\
    Journey Patterns             & 12\,722  \\
    Points in Journey Pattern    & 461\,359 \\
    Service Journeys             & 12\,722  \\
    Passing Times                & 461\,359 \\
    Shape Points                 & 91\,850  \\
    \hline
  \end{tabular}
\end{table}
 
Structural fidelity was inspected beyond raw counts. The 71 lines decompose into
138 routes (71 outbound, 67 inbound), with 67 lines carrying both directions and 4
carrying a single directional route. Each route decomposes into journey patterns at an
average of 92.2 per route (min.\,15, max.\,222), reflecting the variance in
stopping patterns across the network, and the mean stop count per pattern is 36.3
(min.\,7, max.\,66). All scheduled stop points appear in at least one pattern, and
each trips.txt record maps to exactly one journey pattern and one service
journey, which yields the 1:1 correspondence
($\lvert J \rvert = \lvert \mathit{JP}\rvert$) established in
Equation~\eqref{eq:fml-trip}.
 
Three day types were derived from the feed calendar (weekdays, Saturdays, and
Sundays/holidays), together with 365 operating-day exceptions from
calendar\_dates.txt that encode date-level overrides. Service journeys
are distributed as 5\,644 on weekdays, 3\,847 on Saturdays, and 3\,231 on Sundays and
holidays, spanning a service window from 00:00 to 23:55.
 
Figure~\ref{fig:shape_scatter} plots the 91\,850 imported shape points directly
from their persisted longitude and latitude coordinates. The rendered geometry
reproduces the topology of the Porto metropolitan network from the relational store
alone, which confirms that the spatial trajectories survive the GTFS-to-Transmodel
mapping and relational persistence without distortion.
 
\begin{figure}[htbp]
    \centering
    \includegraphics[width=\linewidth]{assets/shape_scatter.png}
    \caption{Scatter plot of the 91\,850 imported shape points, drawn from their
    persisted coordinates with no recourse to the original feed geometry. The
    network topology is reconstructed entirely from the Transmodel relational
    store.}
    \label{fig:shape_scatter}
\end{figure}

\section{Algorithmic Validation and Feed Normalization}
\label{sec:res-validation}

A round-trip transformation validated the export pipeline: the source GTFS feed
was imported, persisted in the Transmodel database, and re-serialized to both GTFS
and NeTEx. Each output was validated against its standard: the MobilityData GTFS
Validator v8.0.0~\citep{mobilitydata2026gtfsvalidator} for GTFS, and XSD validation
against the NeTEx CEN v1.3 schema~\citep{netexcen2026xsd} for NeTEx. Both
artifacts passed with zero errors. The re-exported GTFS feed carries 3\,910
warnings against the original feed's 5\,599, a reduction of 30.2\% attributable to
the normalization applied by the Transmodel intermediary layer.

Table~\ref{tab:warning_breakdown} attributes the reduction by warning code. Almost
the entire reduction corresponds to \codeid{mixed_case_recommended_field}, which the
Pydantic-validated Transmodel layer normalizes to consistent casing. The residual
warnings are not pipeline artifacts: \codeid{non_ascii_or_non_printable_char}
reflects the accented characters inherent to a Portuguese-language feed, and the
shape-geometry warnings are present in the source feed and preserved unchanged. The
three \codeid{service_has_no_active_day_of_the_week} warnings introduced by the
round trip correspond to placeholder day types synthesized for
\codeid{calendar_dates}-only services, whose weekday flags are false by
construction.

\begin{table}[htbp]
\centering
\caption{GTFS warning breakdown: original feed versus exported feed.}
\label{tab:warning_breakdown}
\begin{tabular}{lrr}
\hline
\textbf{Warning Code} & \textbf{Original} & \textbf{Exported} \\
\hline
\codeid{mixed_case_recommended_field}         & 1\,696 &     4 \\
\codeid{non_ascii_or_non_printable_char}     & 3\,899 & 3\,900  \\
\codeid{stop_too_far_from_shape}              &      2 &     2  \\
\codeid{stop_too_far_from_shape_using_user_distance} & 1 & 0 \\
\codeid{stops_match_shape_out_of_order}      &      1 &     1  \\
\codeid{service_has_no_active_day_of_the_week} & -- &  3  \\
\hline
\textbf{Total}                                    & 5\,599 & 3\,910 \\
\hline
\end{tabular}
\end{table}

Schema validation governs the NeTEx publication chain. National Access Points
reject documents that fail XSD validation before any semantic check is applied.
Passing the CEN schema, which spans thousands of element definitions, is therefore
the minimum machine-checkable evidence that an export produces documents downstream
systems can accept and parse. The exported NeTEx XML passed XSD validation against
the NeTEx CEN v1.3 schema with zero errors. The resulting Publication Delivery
document organizes the network across five frames: a ResourceFrame (the operator),
a SiteFrame (2\,550 StopPlaces), a ServiceCalendarFrame (3 DayTypes), a
ServiceFrame (71 Lines, 2\,550 ScheduledStopPoints, and the 220 deduplicated
JourneyPatterns), and a TimetableFrame (12\,722 ServiceJourneys with 461\,359
TimetabledPassingTimes).

The exporter resolves the GTFS-to-NeTEx time-representation gap: 9\,486 passing
times in the case-study feed use hours $\geq 24$, which is valid in GTFS for
overnight service but invalid under \texttt{xs:time}. These are normalized to valid
\texttt{xs:time} values with an accompanying day-offset element (24:10:00 becomes
00:10:00 with a \texttt{DepartureDayOffset} of 1).
 
\section{Computational Performance and Data Compression}
\label{sec:res-performance}
 
Figure~\ref{fig:pipeline_dashboard} reports the artifacts produced or consumed
during pipeline execution (panel a) and the wall-clock duration of each step
(panel b). The SQLite database (104.4\,MB) is approximately 20$\times$ larger
than the compressed GTFS input (5.1\,MB), reflecting full relational indexing
with foreign keys across all Transmodel entities and the lossless retention of
unmapped source fields as per-row \texttt{extras} records, whereas the exported
GTFS archive (5.0\,MB) is marginally smaller than the original because both
include shapes and the CSV serialization normalizes redundant formatting.
 
\begin{figure}[htbp]
    \centering
    \includegraphics[width=\textwidth]{assets/fig_pipeline_dashboard.png}
    \caption{Pipeline artifacts and step durations on the Porto feed. Panel (a)
    reports the size of each artifact produced or consumed; the SQLite store is a
    temporary artifact created only for the validation benchmark. Panel (b) reports
    per-step wall-clock duration on a single-threaded process (log scale), colored
    by stage. NeTEx XSD validation is an optional conformance check rather than a
    required export step.}
    \label{fig:pipeline_dashboard}
\end{figure}
 
The NeTEx output is the largest artifact, yet it is smaller than a direct
one-pattern-per-trip serialization would produce: the journey-pattern deduplication
pass of \S\ref{sec:fml-dedup} collapses the 12\,722 structurally identical patterns
into 220 canonical representatives (a 58$\times$ reduction in pattern count),
eliminating repeated \texttt{pointsInSequence} elements and their nested
\texttt{StopPointInJourneyPattern} children; the residual size is dominated by the
461\,359 timetabled passing times, each serialized as a multi-element XML subtree.
 
The duration profile (Figure~\ref{fig:pipeline_dashboard}, panel~b) shows that
GTFS import (32.1\,s) and GTFS export (12.7\,s) remain interactive, so an
operator can ingest a feed, correct an error, and re-export within a single
working session. NeTEx export (45.7\,s) is the dominant generation
step, reflecting the construction and serialization of over 1.4 million XML
elements and the deduplication pass that computes stop-sequence signatures across
all 12\,722 patterns. NeTEx XSD validation is the longest step overall, since
\texttt{lxml} validates each element against the nested NeTEx schema hierarchy; it
is an optional conformance check rather than a required export step. The
import--export stages of the round-trip workflow complete in approximately 91\,s in
aggregate on a single-threaded process (means over repeated isolated runs),
which bounds the wall-clock cost of importing a feed into the Transmodel
relational schema and exporting both GTFS and NeTEx representations. All steps
are CPU-bound: CPU occupancy stays between 78\% and 97\% of wall-clock time
across the pipeline.
 
\section{National-Scale NeTEx Ingestion: The Swiss Federal Feed}
\label{sec:res-netex-scale}

The Porto case study validates the GTFS ingestion path and the round-trip export
chain, but its NeTEx leg consumes only documents that PID-transit itself produced.
To evaluate the NeTEx importer against data authored by an independent publisher,
the toolkit was applied to the Swiss national timetable for 2026, distributed
through the national access point under the Swiss SKI profile: 355 XML files
totalling 24.5\,GB uncompressed (a 420\,MB archive). The dataset exercises
profile constructs absent from both the Porto feed and the toolkit's own dialect:
entities are referenced across file boundaries in arbitrary order (the stop
registry, calendar, and per-line timetable files are separate documents),
journeys carry no journey patterns and instead embed their stop sequence as
inline \texttt{Call} elements, and calendars are encoded as
\texttt{AvailabilityCondition} validity bitmasks rather than day-of-week
properties.

The import completed without manual preprocessing and yielded 482 operators,
4\,145 lines, 67\,183 scheduled stop points (every one resolved to coordinates
through the \texttt{PassengerStopAssignment}--\texttt{Quay} indirection),
69\,652 day types with 3\,717\,297 operating-day exceptions materialized from
the validity bitmasks, approximately 1.66 million service journeys, and
approximately 26.5 million timetabled passing times, persisted as an 8.64\,GB
relational store. The run took 108.4 minutes of wall-clock time, of which 45.9
minutes were CPU time: at this scale the pipeline becomes I/O-bound on the
database volume (CPU occupancy of 42\%), whereas on feeds whose database fits
the page cache it is CPU-bound (occupancy of 78--97\%). Peak resident memory was
1.17\,GB, and, decisively, this figure is identical to that of a 630\,MB
single-timetable subset of the same feed: the two-pass streaming importer
(\texttt{iterparse} with element pruning) holds memory flat regardless of input
size, so national-scale ingestion is bounded by disk throughput rather than by
RAM, and completes on commodity hardware.

Figure~\ref{fig:netex_network} repeats the geometric-fidelity argument of
Figure~\ref{fig:shape_scatter} for the Swiss import. The SKI profile ships no
shape geometry (\texttt{LinkSequenceProjection} is absent), so the network is
reconstructed from the relational model itself: the left panel scatters the
67\,183 imported stop coordinates, and the right panel draws 30\,477 distinct
edges connecting consecutive stops of each journey's passing-time sequence,
computed with a window function over the \texttt{passing\_time} table for a
sample of 50 of the 348 per-line timetable files (529 lines). The recognizable
Swiss network topology --- the Mittelland corridor, the principal urban
clusters, and the Alpine valley lines --- emerges from persisted coordinates and
sequence ordering alone, confirming that both survive the cross-file NeTEx
import intact.

\begin{figure}[htbp]
    \centering
    \includegraphics[width=\textwidth]{assets/netex_network_scatter.png}
    \caption{Swiss national network reconstructed solely from the Transmodel
    relational store after NeTEx import. Left: raw scatter of the 67\,183
    scheduled stop points (1\,222 international outliers lie outside the frame).
    Right: 30\,477 distinct consecutive-stop edges derived from
    \texttt{passing\_time} sequence order, colored by line, for a sample of 50
    of the 348 per-line timetable files (529 lines). The SKI profile publishes
    no shape geometry, so stop coordinates and passing-time ordering are the
    complete geometric content of the feed; no basemap is used.}
    \label{fig:netex_network}
\end{figure}

\section{Comparative Architectural Benchmarking}
\label{sec:res-benchmark}
 
This section compares PID-transit against one representative of each architectural
class: a file-to-file converter (the
Theoremus Go converter) and a persistent system of record (chouette-core).
 
\subsection*{Experimental conditions}
 
The two head-to-head comparisons were executed on different hosts: chouette-core
requires a Docker-based PostGIS and Rails stack available only on a Linux
workstation, whereas the Theoremus comparison ran on the Windows benchmarking host.
To keep each pairing internally valid, PID-transit
was re-measured on the same host as the respective competitor, so
Figure~\ref{fig:benchmark} reports two same-host groupings that are not
cross-comparable. The file-to-file pairing reflects the current release of
PID-transit (including the streaming NeTEx importer); the system-of-record
pairing retains the measurements taken on the Linux workstation with the
release available at that time. Two quantities are reported: wall-clock execution time, taken as
the mean over seven repetitions with one warm-up discarded, and peak resident set
size (peak RSS), measured with no container memory limit. For chouette-core,
worker-process RSS is reported separately from the standing PostGIS server.
 
\subsection{Profile Interoperability}
\label{sec:res-bench-profiles}
 
A shared-input timing comparison is not directly achievable because each tool
consumes only a subset of NeTEx profile families. The EPIP export resolves stops
through constructs that the Nordic-profile Theoremus loader does not read, and a
European-profile feed exercises \texttt{LineRef} indirections that the EPIP importer
does not expose; in both directions the feed parses but the structural payload is
dropped. Strict profile interoperability is therefore a primary practical
differentiator among converters, and throughput becomes relevant only after profile
compatibility is established. To obtain an equal-work comparison, a benchmarking
compatibility shim was applied to the NeTEx exporter (emission of top-level
\texttt{<Quay>} elements and a \texttt{QuayRef} attribute on each
\texttt{ScheduledStopPoint}). The shim affects only profile addressing, not EPIP
compliance, and the round trip remains lossless: the 461\,359 passing times are
preserved and the unit-test suite passes.
 
\subsection{Quantitative Comparison}
\label{sec:res-bench-quant}
 
With the shim applied, Theoremus converted the Porto export completely (2\,550
stops, a 31\,MB stop\_times.txt), which establishes an equal-work comparison.
For chouette-core, the import produced 461\,359 passing times and 12\,722 vehicle
journeys identical to those of PID-transit and deduplicated to the same 220 journey
patterns (\S\ref{sec:fml-dedup}). Figure~\ref{fig:benchmark} reports the two
same-host pairings, contrasting execution time (panel a) and peak memory (panel b).
 
\begin{figure}[htbp]
    \centering
    \includegraphics[width=\textwidth]{assets/fig_benchmark.png}
    \caption{Same-host head-to-head comparison on the Porto feed (2\,550 stops,
    12\,722 journeys, 461\,359 passing times): (a) execution time (log scale) and
    (b) peak resident memory. Bars pair PID-transit with one competitor per
    operation; error bars denote one standard deviation over seven repetitions (one
    warm-up discarded). The dotted divider separates the file-to-file pairing (vs
    Theoremus, Go) from the system-of-record pairing (vs chouette-core); the two
    groupings were measured on different hosts and are not cross-comparable.}
    \label{fig:benchmark}
\end{figure}
 
Against the compiled, streaming Go converter (Theoremus), PID-transit completes
the same NeTEx-to-GTFS conversion in 30.5\,s against 17.5\,s, a factor of about
1.7, executed as two stages (NeTEx import, $19.1 \pm 0.4$\,s; GTFS export,
$11.4 \pm 1.3$\,s) whose intermediate product is the queryable relational store
that the file-to-file converter does not retain. The memory comparison favors
PID-transit where scale matters: its streaming NeTEx import
(\texttt{iterparse} with element pruning, the same mechanism that enables the
national-scale ingestion of Section~\ref{sec:res-netex-scale}) peaks at
80\,MB, roughly 6.7$\times$ below the converter's 539\,MB single-step
footprint, while its own largest stage (GTFS export, 594\,MB) is comparable to
that footprint. Run-to-run variation of the import stage stays near 2\%,
against roughly 18\% for the converter.
 
Against the system-of-record peer (chouette-core), PID-transit produces
byte-equivalent output between 9$\times$ (export) and 25$\times$ (import) faster;
normalized to feed size, it imports at 2.8\,s per 100\,000 passing times against
70\,s for chouette-core. These ratios reflect the cost of platform generality
rather than implementation efficiency alone, since chouette-core additionally
performs PostGIS spatial storage, per-import schema isolation, and dataset
versioning. The memory result is mixed (Figure~\ref{fig:benchmark}, panel~b):
chouette-core uses less RAM on export because it streams one NeTEx file per line,
whereas PID-transit builds the whole XML tree in memory, its known bottleneck
(\S\ref{sec:res-performance}), and additionally avoids the standing PostGIS server
that chouette-core maintains. Deployment-footprint benchmarking details are reported
in the supplementary material under \emph{Footprint benchmarking details}.
 
\section{Practical Implications}
\label{sec:implications}
 
The computational results establish standards compliance and feasibility; their
translation into practice depends equally on the institutional capacity of the
operators expected to adopt these standards. Two implications follow.
 
The deployment shows that regulatory compliance under EU Delegated Regulation
2024/490 is reachable without an enterprise platform: a single-process toolkit
imports an operator's existing feed and exports both GTFS and NeTEx conforming to
EPIP, so the engineering cost of National Access Point submission is bounded rather
than open-ended. The standardized relational layer also positions an operator's
schedule data as a reusable input for the energy-and-mobility applications that
decarbonization programs require. Electric-bus depot scheduling, charging-window
planning, and the coordination of fleet charging with the constraints communicated
by distribution system operators (DSOs) depend on accurate service spans,
vehicle-hours, and journey patterns, the quantities the analytics layer of
\S\ref{sec:fml-metrics} derives directly from $T$. Because every exchange format
is treated as an encoding of the single canonical Transmodel model
(\S\ref{sec:fml-translation}), a planned extension that adds a dedicated GTFS
table for electric-vehicle charging infrastructure
propagates to the NeTEx export without a separate data model, so charging assets
can be described in the same feed as the service they support. Regulatory
compliance thus constitutes a practical engineering task rather than a purely
administrative obstacle, and a single standardization effort serves both the
immediate MMTIS reporting obligation and the subsequent integration of
electric-fleet operations into the wider energy system.
 
The technical maturity of GTFS, Transmodel, and NeTEx is not the principal barrier
to adoption; the scarcity of training material outside English and French is. A
Portuguese-language curriculum is released as an open
repository (available on request)
under a permissive license so that Portuguese-speaking operators, instructors, and
regulators can reuse and adapt it. Each module follows a consistent structure
(learning objectives, a conceptual motivation, references to the relevant
specification clauses, an annotated worked example in NeTEx with its GTFS
counterpart where applicable, a brief GTFS--NeTEx mapping note, and practical
exercises with solutions), supporting both linear study and later reference. The
curriculum maps NeTEx concepts to their GTFS files and Transmodel sub-models across
nine modules, progressing from core framing to the Portuguese national profile and
practical tooling; the full mapping is provided in the supplementary material under
\emph{NeTEx training curriculum structure and mapping to GTFS and Transmodel}.
Pairing the toolkit with this curriculum reduces the
educational barrier that conceptual complexity and documentation language would
otherwise impose on small and medium operators seeking standards-conformant data
exchange.

\section*{Declaration of generative AI and AI-assisted technologies in the manuscript preparation process}

During the preparation of this work the authors used Claude (Anthropic) in order
to assist with the development of the software code described in this manuscript.
After using this tool, the authors reviewed and edited the content as needed and
take full responsibility for the content of the published article.


%% Notas:
%% 1. O capítulo de exemplo exibe informações e exemplos sobre comandos e
%%    ambientes TeX/LaTeX, assim como os específicos do pacote utfpr-thesis,
%%    entre outros. Estes comandos e ambientes podem ser utilizados para
%%    produzir as diversas estruturas necessárias à elaboração do documento:
%%    enumerações; citações e referências; equações; ilustrações; tabelas;
%%    abreviaturas, siglas e acrônimos; símbolos; glossário; apêndices e anexos;
%%    índice; entre outras. Observar comentários em todos os arquivos-fonte do
%%    projeto para entender configurações e formatações no modelo.
%% 2. O arquivo final (PDF) pode ser convertido para PDF/A usando diversas
%%    ferramentas, por exemplo:
%%      https://www.pdfforge.org/online/en/pdf-to-pdfa